\documentclass[11pt,a4paper]{article}
\usepackage{ae,lmodern}
\usepackage[utf8]{inputenc}
\usepackage[T1]{fontenc}
\usepackage{amsfonts}
\usepackage[french]{babel}
\frenchbsetup{StandardLists=true} % à inclure si on utilise \usepackage[francais]{babel}
\usepackage{enumitem}
\usepackage{amssymb}
\def\nbR{\ensuremath{\mathrm{I\! R}}}
\usepackage{amsmath}
\DeclareMathOperator{\e}{e}

\usepackage[left=2cm,right=2cm,top=2cm,bottom=2cm]{geometry}
\usepackage{multirow}
\usepackage[dvipsnames]{xcolor}
\usepackage{parskip}
\usepackage{caption}
\usepackage{calc}




\usepackage[T1]{fontenc}
\usepackage{fouriernc}

\usepackage{graphicx}
\usepackage{setspace}
\singlespacing

\usepackage{tcolorbox}
\usepackage{framed}
\usepackage{multicol}

\usepackage{fancybox}
\usepackage{fancyhdr}
\pagestyle{fancy}
\renewcommand\headrulewidth{1pt}
\fancyhead[L]{Lycée Denis Diderot }
\fancyhead[R]{Classe 1BTS}
\setcounter{page}{1}\fancyfoot[C]{Chapitre 6 - Probabilités conditionnelles -  \thepage}
\fancyfoot[R]{2024-25}
\fancyfoot[L]{Alexis RUIZ}
\usepackage{pstricks,pst-plot}
\usepackage{tikz}
\usepackage{tkz-tab}
\usetikzlibrary{arrows}
\usepackage{pifont}

\newcommand{\ud}{\mathrm{d}}
\newcommand{\V}{\overrightarrow}
\newcommand{\Rep}{(O;\V{u};\V{v})} 

\newcounter{numexos}%Création d'un compteur qui s'appelle numexos
\setcounter{numexos}{0}%initialisation du compteur
\newcommand{\exercice}[1]{%Création d'une macro ayant un paramètre
\addtocounter{numexos}{1}%chaque fois que cette macro est appelée, elle ajoute 1 au compteur numexos
\textcolor{black}{\normalsize Exercice\,\thenumexos\,}\,#1%Met en noir Exercice et la valeur du compteur appelée par \thenumeexos
}

\newcounter{numact}%Création d'un compteur qui s'appelle numexos
\setcounter{numact}{0}%initialisation du compteur
\newcommand{\activit}[1]{%Création d'une macro ayant un paramètre
\addtocounter{numact}{1}%chaque fois que cette macro est appelée, elle ajoute 1 au compteur numexos
\textcolor{black}{A\thenumexos\,}\,#1%Met en noir Exercice et la valeur du compteur appelée par \thenumeexos
}

\newcounter{nump}%Création d'un compteur qui s'appelle numexos
\setcounter{nump}{0}%initialisation du compteur
\newcommand{\propriete}[1]{%Création d'une macro ayant un paramètre
\addtocounter{nump}{1}%chaque fois que cette macro est appelée, elle ajoute 1 au compteur numexos
\textcolor{black}{Propriété\,\thenump\,}\,#1%Met en noir Exercice et la valeur du compteur appelée par \thenumeexos
}



\newcommand{\bex}[2]
{\begin{bclogo}[couleur = white,logo= , barre=none,  arrondi = 0.1]
{~\exercice:~~{#1}}
~~
{#2}
~~
\end{bclogo}}

\newcommand{\act}[3]
{\begin{bclogo}[couleur = white,logo=,  arrondi = 0.1,barre=none]
{~\activit:~~{#2}}
~~
{#2}
~~
~~
\end{bclogo}
\vspace{3mm}
\cadre{{#3}}
}

\newcommand{\rem}[1]
{\begin{bclogo}[couleur = white,logo=\bcinfo , arrondi = 0.1]
{~Remarque~:}
~~
{#1}
~~
\end{bclogo}}

\newcommand{\bpre}[1]
{\begin{bclogo}[couleur = white,logo=\bcinfo , arrondi = 0.1]
{~\normalsize Prérequis~:}
~~
{#1}
~~
\end{bclogo}}



\newcommand{\intro}[1]
{\begin{bclogo}[couleur = white,logo=\bcinfo , arrondi = 0.1]
{~Introduction~:}
~~
{#1}
~~
\end{bclogo}}


\newcommand{\bdef}[1]
{\begin{bclogo}[couleur = white,logo=\bcinfo , arrondi = 0.1]
{~Définition~:}
~~
{#1}
~~
\end{bclogo}}


\newcommand{\bprop}[2]
{\begin{bclogo}[couleur =white ,logo= , barre=none,  arrondi = 0.1]
{~\propriete:~~{#1}}
~~
{#2}
~~
\end{bclogo}}




\newcommand{\cadre}[1]
{\begin{bclogo}[couleur = white,logo=,  arrondi = 0.1,barre=none]
~~
{%
\rule{0mm}{#1mm}}\par
\medskip
~~
\end{bclogo}}


\newcommand{\cadreblanc}[1]{%
\medskip\framebox[\linewidth][c]{%
\rule{0mm}{#1mm}}\par
\medskip}

\setlength{\columnseprule}{.25mm}
\usepackage{multirow,tabularx,booktabs}
\usepackage{colortbl}
\usepackage{wrapfig}

\usepackage[tikz]{bclogo} 
\newcolumntype{C}[1]{>{\centering\arraybackslash}p{#1cm}}
%%%%%%%%%%%%%%%%%%%%%%%%%%%%%%%%%%%%%%%%%%%%%%%%%%%%%%%%%%%%%%%%%%%%%%%%%%%%%%
%%%%%%%%%%%%%%%%%%%%%%%%%%%%%%%%%%%%%%%%%%%%%%%%%%%%%%%%%%%%%%%%%%%%%%%%%%%%%%
%%%%%%%%%%%%%%%%%%%%%%%%%%%%%%%%%%%%%%%%%%%%%%%%%%%%%%%%%%%%%%%%%%%%%%%%%%%%%%


\begin{document}
\newtcolorbox{mybox}{colback=white,
colframe=black!5!black}
\begin{mybox}
\begin{center}    
\textbf{\Large CHAPITRE 6 - LES PROBABILITÉS CONDITIONNELLES}\end{center}
\end{mybox}
\vfill

\begin{bclogo}[couleur = white,logo=\bcinfo,  arrondi = 0.1]
{~Ojectifs~:}
~~
\begin{itemize}
\item Construire un arbre ou un tableau de probabilités en lien avec une situation donnée.
\item Explorer l'arbre ou le tableau de probabilités pour déterminer des probabilités.
\item Calculer la probabilité d'un événement connaissant ses probabilités conditionnelles relatives à une partition de l'univers.
\item Ultiser ou justifier l'indépendance de deux événements.  
\end{itemize}
~~
\end{bclogo}

\vfill 


\bex{Rappels }{}
\cadre{1}


\begin{enumerate}
	\item $P(A)=0,7 ~~~~ P(B)=0,5 ~~~~ P(A\cap B)=0,3$. Calculer $P(A\cup B)$.\\[3pt] 
    
	\dotfill
    
	\item $A$ et $B$ sont incompatibles $P(A)=0,2 ~~~~ P(B)=0,2 $. Calculer $P(A\cup B)$.\\[3pt] 
    
	\dotfill
    
	\item $P(\bar{A})=0,2 ~~~~ P(B)=0,4 ~~~~ P(A\cap B)=0,2$. Calculer $P(A\cup B)$.\\[3pt] 
    
	\dotfill
    
\end{enumerate}


\vfill 


\bex{ Jeux à gratter} { 
Une société souhaite mettre sur le marché un nouveau jeu à gratter dont le principe est le suivant : 
   \begin{itemize}
   		\item la première case représente soit une soleil soit une lune.  
   		\item la seconde case représente soit un coeur soit un pique.
   		\item la société souhaite commercialiser 10 000 tickets répondant aux contraintes suivantes : 
   			\begin{itemize}
   			\item 10\% de ces tickets comportent un soleil
   			\item 1\% des tickets présentant un soleil comportent un cœur. 
   			\item 0,5\% des tickets présentant une lune comportent un cœur.  
   			\end{itemize}
   \end{itemize}
   }

  
   	  1. Compléter le tableau suivant : \\
      \begin{center}
          
   	  \begin{tabular} { | c | c | c | c | }
   \hline
 \rowcolor{white}     & Avec un soleil & Avec une lune & Total \\
   \hline 
   \rowcolor{white}     Avec un coeur  & ~~ & ~~ & ~~ \\
   \hline
   \rowcolor{white}      Avec un pique  & ~~ & ~~ & ~~  \\
   \hline
   \rowcolor{white}      Total  & ~~ & ~~ & 10 000   \\
   \hline
 \end{tabular}
  \end{center}
  
 2.  Un joueur achète un ticket. Calculer la probabilité que ce ticket comporte un cœur.
 \cadre{5}


 \vfill

 \bex{QCM}
 {
 \vspace{2mm}
 Deux machines A et B produisent une même type de pièce. On a prélevé 3 000 unités sortant de la machine A et 2 000 de la machine B. Ces pièces peuvent présenter deux types de défauts : un défaut de couleur noté C, et un défaut de taille noté T.\\
Pour la machine A, 2\% des pièces présentent uniquement le défaut C, 5\% uniquement le défaut T et 1\% les deux défauts.\\
Pour la machine B, 3\% des pièces présentent uniquement le défaut C, 4\% uniquement le défaut T et 2\% les deux défauts.\\
On pourra éventuellement se servir du tableau ci-dessous : \\[5 pt]
 

\begin{center}
\begin{tabular} { | c | c | c | c | c | c | }
   \hline
 \rowcolor{white}      & C seul & T seul & C et T & Ni C ni T & Total \\
   \hline 
 \rowcolor{white}   Machine A  & ~~ & ~~ & ~~ & ~~& ~~\\
   \hline
\rowcolor{white}    Machine B  & ~~ & ~~ & ~~& ~~& ~~\\
   \hline
\rowcolor{white}    Total  & ~~ & ~~ & ~~& ~~ & 5 000  \\
   \hline
 \end{tabular}
 \end{center}
 \vspace{4mm}
 
 On prend au hasard une pièce parmi les 5 000 prélevées ; toutes les pièces ont la même chance d'être choisies. 
}



\begin{multicols}{2}

1. la probabilité que la pièce soit fabriquée par la machine A est : 

\begin{itemize}[label=$\square$]
             \item   a) $\dfrac{3}{5}$
  
             \item   b) $\dfrac{2}{5}$
  
             \item   c) Le tableau ne permet pas de conclure
\end{itemize}
\bigskip


2. la probabilité que la pièce présente uniquement le défaut C est  : 

\begin{itemize}[label=$\square$]
             \item   a) $0,024$
  
             \item   b) $0,02$
  
             \item   c) $0,03$
             
             \item   d) $0,120$
             
\end{itemize}
\columnbreak


3. la probabilité que la pièce présente uniquement le défaut T est  : 

\begin{itemize}[label=$\square$]
             \item   a) $\dfrac{23}{500}$
  
             \item   b) $\dfrac{3}{50}$
  
             \item   c) $\dfrac{7}{500}$
             
             \item   d) $\dfrac{1}{20}$
             
\end{itemize}
\bigskip

4. la probabilité que la pièce présente au moins l'un des deux défauts est : 

\begin{itemize}[label=$\square$]
             \item   a) $0,014$
  
             \item   b) $0,06$
  
             \item   c) $0,038$
             
             \item   d) $0,084$
             
\end{itemize}
\bigskip
\end{multicols}
}

\vfill

\bpre{Pièces défectueuses dans la production 

Dans un atelier, deux machines A et B produisent des pièces identiques. 40\% proviennent de A dont 5\% sont défectueuses. D'autre part 2\% des pièces produites par B sont défectueuses.  \\[2 pt] }

1. A l'aide des données de l'énoncé compléter le tableau suivant :\\
\begin{center}
\begin{tabular} { | c | c | c | c | }
   \hline
 \rowcolor{white}     & Produites par A & Produites par B & Total \\
   \hline 
 \rowcolor{white}  Défectueuses  & ~~ & ~~& ~~\\
   \hline
\rowcolor{white}   Non défectueuses   & ~~& ~~& ~~\\
   \hline
\rowcolor{white}   Total  & ~~& ~~ & 5 000  \\
   \hline
 \end{tabular}
\end{center}
 
\vspace{5pt}

2. On prélève au hasard une pièce dans la production. Toutes les pièces sont supposées être identiques et indiscernables au toucher. Les tirages sont donc équiprobables. \\[2 pt]
Calculer la probabilité des évènements suivants :
\begin{multicols}{3}
A : "la pièce à été fabriquée par la machine A"\\[2 pt]
\cadre{1}
B : "la pièce à été fabriquée par la machine B"\\[2 pt]
\cadre{1}
D : "la pièce est défectueuse"\\[2 pt]
\cadre{4}
\end{multicols}


3. Définir par une phrase l'évènement A $\cap$ D puis calculer sa probabilité.\\[2 pt]
\cadre{2}




4. On prélève au hasard une pièce parmi l'ensemble des pièces défectueuses. Calculer la probabilité que la pièce ait été produite par la machine A. \\[2 pt]
\cadre{15}

5. Calculer $\dfrac{P(A \cap D)}{P(D)}$ et comparer avec le résultat précédent.\\[2 pt]
\cadre{5}



\textbf{Remarque : }\\
\cadre{1}
\vspace{0.5cm}



 \vspace{2mm}

\bpre{\textbf{Tirage dans une urne } \\
 \vspace{-2mm}
Une urne contient 3 boules blanches et 2 boules rouges indiscernables au toucher.
On tire successivement deux boules sans remettre la première dans l'urne. \\[2pt]
}
\begin{multicols}{2}
~~\\
\begin{center}

\begin{tikzpicture}[scale=0.5]
\fill (0,0) circle (1pt);
\node[circle,draw] (b11) at (2,2) {R};

\node[circle,draw] (b13) at (2,-2) {B};

\node[circle,draw] (b21) at (5,3) {R};
\node[circle,draw] (b22) at (5,1) {B};
\node[circle,draw] (b23) at (5,-1) {R};
\node[circle,draw] (b24) at (5,-3) {B};



\draw (0,0) -- (b11.south west);

\draw (0,0) -- (b13.north west);

\draw (b11.north east) -- (b21.west);
\draw (b11.south east) -- (b22.west);
\draw (b13.north east) -- (b23.west);
\draw (b13.south east) -- (b24.west);

\end{tikzpicture}
    
\end{center}

1. Déterminer le nombre de tirages distincts possibles.\\[1pt]
\cadreblanc{8}
2. Calculer la probabilité d'obtenir deux boules rouges.\\[1pt]
\cadreblanc{8}
3. Calculer la probabilité d'obtenir deux boules blanches.\\[1pt]
\cadreblanc{8}
4. Calculer la probabilité d'avoir deux boules de couleur différente. \\[1pt]
\cadreblanc{8}

\end{multicols}

\bpre{\textbf{Contrôle de production}

Le contrôle de la production de 10 000 pièces dans un atelier a donné les résultats suivants : }
\begin{center}
\begin{tabular} { | c | c | c | c | }
   \hline
      & Défaut de longueur & Pas de défaut de longueur & Total \\
   \hline 
   Défaut de surface  & 6 & 294 & 300\\
   \hline
   Pas de défaut de surface   &194& 9 506&  9700\\
   \hline
   Total  & 200 & 9800 & 10 000  \\
   \hline
 \end{tabular}
\end{center}

1. On prélève au hasard une pièce dans la production totale . Les tirages sont équiprobables. Calculer la probabilité des événements suivants.
\begin{multicols}{2} 
	L : "la pièce présente un défaut de longueur"\\
	\cadreblanc{8}
	S : "la pièce présente un défaut de surface"\\
	\cadreblanc{8}
\end{multicols}
2. Calculer la probabilité de l'évènement S $\cap$ L. Comparer $P(S \cap L)$ et $P(S) \times P(L)$.\\
\cadre{-2}
\dotfill

%%%%%%%%%%%%%%%%%%%%%%%%%%%%%%%%%%%%%%%%%%%%%%%%%%
\newpage
%%%%%%%%%%%%%%%%%%%%%%%%%%%%%%%%%%%%%%%%%%%%%%%%%%
\intro{
On fait remonter à la correspondance de $1654$ entre Pascal et Fermat, sur un problème de jeu de hasard, l'acte de naissance du calcul des probabilités.

Après trois siècles de recherche, le calcul des probabilités a pu fournir, au début du $XX^{e}$ siècle, les bases théoriques nécessaires au développement de la statistique et a investi de très nombreux domaines de la vie scientifique, économique et sociale.

Le but des probabilités est d'essayer de rationaliser le hasard : quelles sont les chances d'obtenir un résultat suite à une expérience aléatoire ?\\
Quelles chances ai-je d'obtenir "pile" en lançant une pièce de monnaie ? Quelles chances ai-je d'obtenir "$6$" en lançant un dé ? Quelles chances ai-je de valider la grille gagnante du loto ?
}


%%%%%%%%%%%%%%%%%%%%%%%%%%%%%%%%%%%%%%%%%%%%%%%%%

\begin{mybox}
\begin{center}    
\textbf{1. Vocabulaire des événements}
\end{center}
\end{mybox}

\begin{mybox}
\begin{center}    
\textbf{1.1 Vocabulaire}
\end{center}
\end{mybox}



   \begin{dinglist}{228}
      \item Chaque résultat possible et prévisible d'une expérience aléatoire est appelé \underline{éventualité} liée à \\ \hspace*{0.3cm} l'expérience aléatoire.
      \item L'ensemble formé par les éventualités est appelé \underline{univers}, il est très souvent noté $\Omega$.
      \item Un \underline{événement} d'une expérience aléatoire est une partie quelconque de l'univers,
      \item Un événement ne comprenant qu'une seule éventualité est un \underline{événement élémentaire}.
      \item L'événement qui ne contient aucune éventualité est l'\underline{événement impossible}, noté $\varnothing$,
      \item L'événement composé de toutes les éventualités est appelé \underline{événement certain}.
      \item Pour tout événement $A$ il existe un événement noté $\overline{A}$ et appelé \underline{événement contraire} de $A$ \\ \hspace*{0.3cm} qui est composé des éléments de $\Omega$ qui ne sont pas dans A.\\
   \hspace*{0.3cm} On a en particulier $A\cup\overline{A}=\Omega$ et $A\cap\overline{A}=\varnothing$.
   \end{dinglist}


\smallskip
\cadre{70}

Dans toute la suite du cours, on suppose que $\Omega$ est l'univers associé à une expérience aléatoire, et $A$ et $B$ deux événements associés à cet univers.

\newpage 
%%%%%%%%%%%%%%%%%%%%%%%%%%%%%%%%%%%%%%%%%%%%%%%%%%%%%%%%%%%%%%%%%%%%%%%%%%

\begin{mybox}
\begin{center}    
\textbf{1.2. Intersection et réunion des évènements }
\end{center}
\end{mybox}

\fbox{\begin{minipage}{16.6cm}{
   \begin{dinglist}{228}
      \item \underline{Intersection d'événements} : événement constitué des éventualités appartenant à $A$ et à $B$ noté $A\cap B$ (se lit "$A$ inter $B$" ou "$A$ et $B$"),
      \item \underline{Réunion d'événements} : événement constitué des éventualités appartenant à $A$ ou à $B$ noté $A\cup B$ (se lit "$A$ union $B$" ou "$A$ ou $B$").
   \end{dinglist}



  \begin{center}
      
   Si $A\cap B=\varnothing$, on dit que les événements sont \underline{disjoints} ou \underline{incompatibles}.
   \end{center}

}
\end{minipage}
}

\cadre{5}

\cadre{30}

\begin{mybox}
\begin{center}    
\textbf{1.3. Représentation des évènements}
\end{center}
\end{mybox}

\textbf{Diagrammes ou patates}

\begin{center}
    
\begin{multicols}{4}

\begin{tikzpicture}
\newcommand{\cercleA}{(0,0.67) circle (1cm)}
\newcommand{\cercleB}{(0,-0.67) circle (1cm)}
\fill[fill=blue!20,even odd rule] \cercleA \cercleB;
\draw[color=blue] \cercleA;
\draw[color=blue] \cercleB;
\end{tikzpicture}

\cadreblanc{15}

\begin{tikzpicture}
\newcommand{\cercleA}{(0,0.67) circle (1cm)}
\newcommand{\cercleB}{(0,-0.67) circle (1cm)}
\draw[color=blue] \cercleA;
\draw[color=blue] \cercleB;
\begin{scope}
\clip \cercleA;
\draw[color=blue,fill=blue!20] \cercleB;
\end{scope}
\begin{scope}
\clip \cercleB;
\draw[color=blue] \cercleA;
\end{scope}
\end{tikzpicture}

\cadreblanc{15}

\begin{tikzpicture}
\newcommand{\cercleA}{(0,0.67) circle (1cm)}
\newcommand{\cercleB}{(0,-0.67) circle (1cm)}
\begin{scope}
\clip (-1.05,0) rectangle (1.05,1.7);
\draw[color=blue,fill=blue!20] \cercleA;
\end{scope}
\begin{scope}
\clip (-1.05,-1.7) rectangle (1.05,0);
\draw[color=blue,fill=blue!20] \cercleB;
\end{scope}
\end{tikzpicture}

\cadreblanc{15}

\begin{tikzpicture}
\newcommand{\cercleA}{(0,0.67) circle (1cm)}
\newcommand{\cercleB}{(0,-0.67) circle (1cm)}
\fill[fill=blue!20] \cercleB;
\draw[color=blue,fill=white] \cercleA;
\draw[color=blue] \cercleB;
\end{tikzpicture}

\cadreblanc{15}

\end{multicols}
\end{center}

\vfill 

\textbf{Tableaux}

On jette deux dés à quatre faces (tétraèdre régulier) et on calcule le produit des deux nombres obtenus :

\begin{center}
   \begin{tabular}{|c|c|c|c|c|}
      \hline 
      & \textbf{1} & \textbf{2} & \textbf{3} & \textbf{4} \\ 
      \hline
      \textbf{1} & $1$ & $ $ & $ $ & $ $ \\ 
      \hline
      \textbf{2} & $ $ & $ $ & $ $ & $ $ \\ 
      \hline 
      \textbf{3} & $ $ & $ $ & $ $ & $ $ \\ 
      \hline 
      \textbf{4} & $ $& $ $ & $ $ & $ $ \\ 
      \hline 
   \end{tabular} 
\end{center}

\vfill

\newpage

\textbf{Arbres}

On lance une pièce de monnaie trois fois se suite, on peut schématiser cette expérience par un arbre :
\begin{center}
\begin{tikzpicture}[scale=0.5]
\fill (0,0) circle (1pt);
\node[circle,draw] (b11) at (3,4) {P};
\node[circle,draw] (b13) at (3,-4) {F};

\node[circle,draw] (b21) at (6,6) {P};
\node[circle,draw] (b22) at (6,2) {F};
\node[circle,draw] (b23) at (6,-2) {P};
\node[circle,draw] (b24) at (6,-6) {F};

\node[circle,draw] (b31) at (9,7) {P};
\node[circle,draw] (b32) at (9,5) {F};
\node[circle,draw] (b33) at (9,3) {P};
\node[circle,draw] (b34) at (9,1) {F};
\node[circle,draw] (b35) at (9,-1) {P};
\node[circle,draw] (b36) at (9,-3) {F};
\node[circle,draw] (b37) at (9,-5) {P};
\node[circle,draw] (b38) at (9,-7) {F};


\draw (0,0) -- (b11.south west);

\draw (0,0) -- (b13.north west);

\draw (b11.north east) -- (b21.west);
\draw (b11.south east) -- (b22.west);
\draw (b13.north east) -- (b23.west);
\draw (b13.south east) -- (b24.west);

\draw (b21.north east) -- (b31.west);
\draw (b21.south east) -- (b32.west);
\draw (b22.north east) -- (b33.west);
\draw (b22.south east) -- (b34.west);
\draw (b23.north east) -- (b35.west);
\draw (b23.south east) -- (b36.west);
\draw (b24.north east) -- (b37.west);
\draw (b24.south east) -- (b38.west);


\end{tikzpicture}
\end{center}


%%%%%%%%%%%%%%%%%%%%%%%%%%%%%%%%%%%%%%%%%%%%%%%%%%


\begin{mybox}
\begin{center}    
\textbf{2. Calcul de probabilités}
\end{center}
\end{mybox}

\cadre{45}

\vfill 

\rem{

Dans un exercice, pour signifier qu'on est dans une situation d'équiprobabilité on a généralement dans l'énoncé une expression du type :
  \begin{itemize}
     \item on lance un dé \underline{non pipé},
     \item dans une urne, il y a des boules \underline{indiscernables} au toucher,
     \item on rencontre au \underline{hasard} une personne parmi \dots
  \end{itemize}
}

\vfill 

\newpage

\bprop{}{}
\vspace{8mm}
\begin{center}
\fbox{
\begin{minipage}{10cm}
{Soit $A$ et $B$ deux événements, on a les propriétés suivantes :  

\vspace{3mm}


   \begin{dinglist}{217}
      \item $P(\varnothing)=0$. \vspace{3mm}
      
      \item $P(\Omega)=1$.  \vspace{3mm}
      
      \item $0 \leqslant P(A) \leqslant 1$.  \vspace{3mm}
      
      \item $P(\overline{A})=1-P(A)$.  \vspace{3mm}
      
      \item $P(A\cup B)=P(A)+P(B)-P(A\cap B)$.  \vspace{3mm}
      
   \end{dinglist}
}
\end{minipage}
}
\end{center}


\vfill 

\bex{}{
   On considère l'ensemble $E$ des entiers de $1$ à $20$. On choisit l'un de ces nombres au hasard. 
   }
   \vspace{8mm}
   \begin{enumerate}


       
  \item $A$ est l'événement : " le nombre est multiple de $3$ " :\\[2 pt]
  \cadre{1} 
  
   \item $B$ est l'événement : " le nombre est multiple de $2$ " :\\[2 pt]
  \cadre{1}
   
   \item Calcul des probabilités : calculer $P(A)$, $P(\overline{A})=1-P(A)$, $P(B)$, $P(A \cap B)$ et $P(A \cup B)$\\[2pt]
\cadre {35}

   \end{enumerate}


\vfill 
\newpage

%%%%%%%%%%%%%%%%%%%%%%%%%%%%%%%%%%%%%%%%%%%%%%%%%%%

\begin{mybox}
\begin{center}    
\textbf{3. Probabilités conditionnelles}

\end{center}
\end{mybox}

\vspace{5mm}

\begin{mybox}
\begin{center}    
\textbf{3.1. Définition}
\end{center}
\end{mybox}



\cadreblanc{45}




   \begin{center}
       
  \textbf{On trouve aussi la notation $P(A/B)$ pour $P_B(A)$.}
 \end{center}



\bex{}{
   On considère l'expérience aléatoire consistant à lancer un dé à $6$ faces, équilibré. On suppose que toutes les faces sont équiprobables, et on définit les événements :
   \begin{itemize}
      \item $B$ : "la face obtenue porte un numéro pair" ;
      \item $A$ : "la face obtenue porte un numéro multiple de $3$".
   \end{itemize}
   }

Déterminons la probabilité d'obtenir un numéro multiple de $3$, sachant qu'on a un numéro pair de deux manières différentes.
\begin{multicols}{2}
\cadre{20}
\cadre{20}
\end{multicols}



 
\vspace{-3mm}
\bprop{}{}

\cadre{15}



\fbox{Démonstration :} Pour $P(B)\neq 0$ et $P(A)\neq 0$, on peut écrire : \\
\begin{itemize}
   \item $P_B(A)=\dfrac{P(A\cap B)}{P(B)}$ d'où $P(A\cap B)=P(B)P_B(A).$
   \item $P_A(B)=\dfrac{P(A\cap B)}{P(A)}$ d'où $P(A\cap B)=P(A)P_A(B).$
\end{itemize}

\newpage 

\bprop{}{}
\vspace{5mm}
\cadre{60}


\smallskip


   Ce théorème revient à dire qu'une probabilité conditionnelle relative à un événement $S$ a toutes les propriétés habituelles du calcul des probabilités.


\smallskip

\bprop{}{}
\vspace{3mm}
\cadre{40}



\underline{Démonstration :} \\[2mm]
$P(A)=P(A\cap B)+P(A\cap \overline{B})=P_B(A)~P(B)+P_{\overline{B}}(A)~P(\overline{B})$.

\vfill 


\newpage

\begin{mybox}
\begin{center}    
\textbf{3.2. Exemple de calcul de probabilités conditionnelles par différentes méthodes }
\end{center}
\end{mybox}

\bex{}{
   Dans un atelier, deux machines $M_1$ et $M_2$ découpent des pièces métalliques identiques. $M_1$ fournit $60\%$ de la production (parmi lesquelles $6,3\%$ sont défectueuses), le reste étant fourni par $M_2$ (dont $4\%$ de la production est défectueuse).\\
   La production du jour est constituée des pièces produites par les deux machines, et on en tire en fin de soirée une pièce au hasard (tous les prélèvements sont supposés équiprobables).}
   \begin{enumerate}
      \item Utilisation des formules des probabilités conditionnelles.
      \begin{enumerate}
          \item Quelle est la probabilité de prélever une pièce défectueuse, sachant qu'elle est produite par $M_1$ ?
          
           
          \dotfill 
          \item Quelle est la probabilité de prélever une pièce défectueuse, sachant qu'elle est produite par $M_2$ ?
        
           
          \dotfill 
          \item Quelle est la probabilité de prélever une pièce défectueuse ?
           
          \dotfill 
       \end{enumerate}
       \item Utilisation d'un tableau : 
       on suppose maintenant que la production est composée de $10000$ pièces.
       \begin{enumerate}
          \item Compléter le tableau suivant qui décrit la production du jour : \\
             \begin{tabular}{|l|l|l|l|}
                 \hline
                   & Nombre de pièces    & Nombre de pièces & \\
                
                & produites par $M_1$   & produites par $M_2$ & ~~~~~~~Total ~~~~~~~\\
                 \hline
                Nombre de pièces défectueuses & & & \\
                 \hline
                Nombre de pièces conformes & & & \\
                 \hline
                Total & & & \\
                 \hline
             \end{tabular}
             \vspace{1mm}

         \item Quelle est la probabilité de prélever une pièce défectueuse ?

         \dotfill 
      \end{enumerate}
      \item Utilisation d'un arbre des probabilités conditionnelles
      \begin{enumerate}
         \item Dresser un arbre des probabilités conditionnelles relatif à la situation proposée.
         \begin{center}
             \begin{minipage}{7cm}
                 \cadre{40}
             \end{minipage}
         \end{center}
         \item Quelle est la probabilité de prélever une pièce défectueuse ?
          \vspace{1mm}
          
         \dotfill
      \end{enumerate}
   \end{enumerate}





\newpage

\begin{mybox}
\begin{center}    
\textbf{3.3. Événements indépendants}
\end{center}
\end{mybox}


\cadre{30}
\bex{}{
   On considère le tirage au hasard d'une carte d'un jeu de $32$ cartes.\\
   
   ~~~~~~~~~~~~~~~~~~$A=$ Tirer un as, \hfill $B=$ Tirer un coeur \hfill et  \hfill $C=$ Tirer un as rouge.~~~~~~~~~~~~~~~~~~
      }
   \vspace{3mm}
   Indépendance de $A$ et $B$ :
\cadre{25}   
 \vspace{3mm}
   Indépendance de $B$ et $C$ :
   
\cadre{25}




\smallskip

\rem{
   Dans le cas où $A$ et $B$ sont des événements de probabilités non nulles, on a 
   \begin{itemize}
      \item $P(A\cap B)=P(B)P_B(A)=P(A)P_A(B)=P(A).P(B)$, d'où :
      \item $P_B(A)=P(A)$ et $P_A(B)=P(B)$.
   \end{itemize}
}
\bprop{}{
   Si $A$ et $B$ sont des événements indépendants, alors : $A$ et $\overline{B}$ ; $\overline{A}$ et $B$ ; $\overline{A}$ et $\overline{B}$ sont également des événements indépendants.}

   %%%%%%%%%%%%%%%%%%%%%%%%%%%%%%%%%%
   \newpage
   %%%%%%%%%%%%%%%%%%%%%%%%%%%%%%%%%%%%%



\bex{}{
\begin{center}
\textbf{Dans cet exercice, sauf mention contraire, les résultats approchés sont à arrondir à}~ \boldmath  $10^{-5}$\unboldmath \end{center}

Un atelier d'une usine d'automobiles est chargé de l'assemblage d'un moteur. 
On s'intéresse, un jour donné, à une machine assurant l'installation de la poulie. Cette machine peut connaître une défaillance susceptible d'être détectée par un système d'alerte. 

Le système d'alerte peut aussi se déclencher sans raison. 

Une étude du service de contrôle a montré que :
\begin{enumerate}
\item[\textbullet] La probabilité que la machine soit défaillante est égale à 0,001.
\item[\textbullet] Lorsque la machine est défaillante, l'alarme se déclenche dans 99 cas sur 100.
\item[\textbullet] Lorsque la machine est n'est pas défaillante, l'alarme se déclenche dans 5 sur 1000.
\end{enumerate} 

On note $D$ l'événement : "la machine est défaillante" et on note $A$ l'événement : "l'alerte est donnée". }

\begin{enumerate}
\item  Donner les probabilités $P(D) ~;~ P_{D}(A)$ et $P_{\overline{D}}(A)$.

\begin{multicols}{3}
	\cadre{5}
	\cadre{5}
	\cadre{5}
\end{multicols}

\item \begin{enumerate}
\item Calculer les probabilités des événements suivants $A \cap D$ et $A \cap \overline{D}$\\
\vspace{-1mm}

\cadre{1}
\cadre{1}
\item En déduire $P(A)$\\
\vspace{-1mm}
\cadre{1}
\end{enumerate}
\item   L'alerte est donnée. Calculer la probabilité qu'il s'agisse d'une "fausse alerte ", c'est-à-dire $P_A(\overline{D})$. Arrondir à $10^{-2}$.

\cadre{10}
\end{enumerate}


\bigskip

\newpage

\bex{}{
Une entreprise vend des calculatrices d'une certaine marque. \\
Le service après-vente s'est aperçu qu'elles pouvaient présenter deux types de défaut, l'un lié au clavier et l'autre lié à l'affichage. \\
Des études statistiques ont permis à l'entreprise d'utiliser la modélisation suivante. 
\begin{itemize}
    \item La probabilité pour une calculatrice tirée au hasard de présenter un défaut de clavier est égale à 0,04. 
    \item En présence du défaut de clavier, la probabilité que la calculatrice soit en panne d'affichage est de 0,03. 
    \item En l'absence de défaut de clavier, la probabilité de ne pas présenter de défaut d'affichage est de 0,94. 
\end{itemize}
On note C l'événement "la calculatrice présente un défaut de clavier " et A l'événement "la calculatrice présente un défaut d'affichage " \\
Dans cet exercice, les probabilités seront écrites sous forme de nombres décimaux arrondis au millième. 
}
\begin{enumerate}
    \item Préciser à l'aide de l'énoncé les probabilités suivantes 
    \begin{enumerate}
        \item $\mathrm{P_{\overline{C}}(\overline{A})}$ \dotfill 
        \vspace{3mm}
        \item $\mathrm{P_{\overline{C}}(A)}$ \dotfill 
        \vspace{3mm}
        \item $\mathrm{P(C)}$ \dotfill 
        \vspace{3mm}
    \end{enumerate}
    \item Construire un arbre pondéré décrivant l'énoncé 
    \begin{center}

    \begin{minipage}{10cm}
     \cadre{30}
     \end{minipage}
             
    \end{center}

    \item On choisit une calculatrice de cette marque au hasard 

    \begin{enumerate}
        \item Calculer la probabilité pour que la calculatrice présente les deux défauts. 

        \dotfill 

        \item Calculer la probabilité pour que la calculatrice présente le défaut d'affichage mais pas le défaut de clavier. 

        \dotfill 

        \item En déduire P(A) 

        \dotfill 

        \item Montrer que la probabilité de l'événement "la calculatrice ne présente aucun défaut " arrondie au millième est égale à 0,902. 

        \dotfill 
        
        
    \end{enumerate}


\end{enumerate}

\bex{}{
Deux machines A et B fabriquent des disques. La machine A produit 1500 disques par jours ; la machine B produit 3000 disques par jour. \\
La probabilité pour qu'un disque ait un défaut est de 0,02 sachant qu'il est produit par la machine A et de 0,035 sachant qu'il est produit par la machine B. On tire au hasard un disque dans la production du jour. 
}

\begin{enumerate}
    \item Calculer la probabilité des événements suivants : 
    \begin{enumerate}
        \item A : "Le disque est produit par la machine A" : \dotfill 
        \vspace{3mm}
        \item B : "Le disque est produit par la machine B" : \dotfill 
        \vspace{3mm}
        \item D : "Le disque a un défaut" : \dotfill 
    \end{enumerate}
    \item Le disque a un défaut .
    \begin{enumerate}
        \item Quelle est la probabilité qu'il ait été produit par la machine A ? \dotfill 
        \vspace{3mm}
        \item Quelle est la probabilité qu'il ait été produit par la machine B ? \dotfill 
    \end{enumerate}

\end{enumerate}

\vspace{3mm}

\bex{}{
On arrondira les probabilité au millième. \\
Dans un lycée, le foyer des lycéens a dénombré les élèves utilisant l'application Toktok. La répartition de ces élèves est donnée dans le tableau suivant : 
\begin{center}

        \begin{tabular}{|c|c|c|c|}
        \hline
             & Filles   & Garçons  & Total  \\
        \hline
            Utilisent Toktok & 148 & 171 & 319 \\
        \hline
             N'utilisent pas Toktok & 81 & 50 & 131 \\
        \hline
             Total & 229 & 221  & 450 \\
        \hline
        \end{tabular}

\end{center}
On prélève au hasard une fiche dans le fichier des élèves du lycée. On admettra que toutes les fiches ont la même probabilité d'être prélevées. On note l’événement \textbf{G}, "la fiche prélevée est celle d'un garçon ";  l’événement \textbf{T}, "la fiche prélevée est celle d'un élève qui utilise Toktok "  
}

\begin{enumerate}
    \item Calculer la probabilité de prélever la fiche d'un garçon

    \vspace{3mm}
    \dotfill 

    \item Montrer que la probabilité de prélever la fiche d'un garçon utilisant l'application Toktok est égale à 0,38 . 

    \vspace{3mm}
    \dotfill
    
    \item Calculer la probabilité de prélever la fiche d'une fille sachant que l'élève correspondant n'utilise pas l'application Toktok. 

    
    \vspace{3mm}
    \dotfill

    \item Calculer la probabilité $\mathrm{P_T(G)}$ et interpréter le résultat. 

   
    \vspace{3mm}
    \dotfill 

    
\end{enumerate}

\newpage

\bex{}{
Une société de produit pharmaceutiques fabrique en très grande quantité un type de comprimés. \\
Un comprimé est conforme si sa masse exprimée en gramme appartient à l'intervalle [1,2 ; 1,3]. La probabilité qu'un comprimé soit conforme est de 0,98. \\
On note :\\[-8mm]
\begin{itemize}
    \item A - L'événement un comprimé est conforme 
    \item B - l'événement un comprimé est refusé
    \item un comprimé conforme est accepté avec une probabilité de 0,98
    \item un comprimé qui n'est pas conforme est refusé avec une probabilité de 0,99.
\end{itemize}

}

\begin{enumerate}
 \item Calculer les probabilités suivantes :
 \begin{enumerate}
     \item $\mathrm{P_A(B)}=$ \dotfill 
     \vspace{3mm}
     \item $\mathrm{P(B \cap A)}=$ \dotfill 
     \vspace{3mm}
     \item $\mathrm{P(B \cap \overline{A})}=$ \dotfill 
 \end{enumerate}

 \item Calculer

  \begin{enumerate}
      \item la probabilité qu'un comprimé soit refusé : \dotfill 

      \vspace{3mm}

      \item la probabilité qu'un comprimé soit conforme, sachant qu'il est refusé, à $10^{-4}$

       \vspace{1mm}
       \dotfill 
       
      
  \end{enumerate}

\end{enumerate}


\bex{}{
Au cours d’une journée, un commercial se déplace pour visiter deux de ses clients afin de leur proposer l’achat d’un produit
de grande consommation.\\
Au vu de son expérience, le commercial estime que : \\
– la probabilité que le premier client visité achète le produit est égale à 0,25 ; \\
– si le premier client achète le produit, la probabilité que le deuxième client visité achète le produit est égale à 0,4 ; \\
– si le premier client n’achète pas le produit, la probabilité que le deuxième client visité achète le produit est égale à 0,25 ; \\
On note A l’événement « Le premier client achète le produit ». \\ 
On note B l’événement « Le second client achète le produit ». \\
}

1. Réaliser un arbre pondéré décrivant cette situation.
\begin{center}

    \begin{minipage}{7cm}
     \cadre{20}
     \end{minipage}
             
    \end{center}

2. Calculer la probabilité de l’événement B. \dotfill 

\vspace{2mm}
3. Déterminer la probabilité qu’un seul client achète le produit. 
\dotfill 







\end{document}